\section{Potenzen und Wurzeln}
\label{sec:fundamentals_powers_and_roots}

Potenzen fassen wiederholte Multiplikationen zusammen. Wurzeln kehren das
Potenzieren um. Aus diesem Zusammenhang ergeben sich alle Rechenregeln dieses
Abschnitts.



\subsection{Potenzen}
\label{subsec:fundamentals_powers}

\begin{DefinitionBox}{Potenz}{fundamentals_power}
  Eine Potenz $a^n$ besteht aus der \emph{Basis} $a$ und dem \emph{Exponenten} $n$.
  Der Exponent gibt an, wie oft die Basis mit sich selbst multipliziert wird:
  \begin{minipage}[c]{0.64\linewidth}
    \[
      a^n=
      \begin{cases}
        1, & n=0,\quad a\neq0,\\[0.3em]
        \underbrace{a\cdot a\cdots a}_{n\text{-mal}},
          & n\in\N,\quad n\geq1.
      \end{cases}
    \]
  \end{minipage}\hfill
  \begin{minipage}[c]{0.27\linewidth}
    \centering
    \resizebox{\linewidth}{!}{%
      \begin{tikzpicture}[font=\small,>=stealth]
  \node (power) at (0,0) {\Huge $a^n$};
  \node[anchor=east,text=THMBlue] (base) at (-1.25,-0.9)
    {\Lang{Basis}{base}};
  \node[anchor=west,text=THMBlue] (exponent) at (1.25,0.9)
    {\Lang{Exponent}{exponent}};
  \draw[->,THMBlue,line width=0.7pt] (base.north east) -- (-0.28,-0.22);
  \draw[->,THMBlue,line width=0.7pt] (exponent.south west) -- (0.28,0.2);
\end{tikzpicture}
%
    }
  \end{minipage}
\end{DefinitionBox}



\subsubsection{Potenzgesetze}
\label{subsubsec:fundamentals_power_laws}

Die drei Potenzgesetze vereinfachen Produkte, Quotienten und Potenzen von
Potenzen. Entscheidend ist jeweils, was übereinstimmt.

\begin{TheoremBox}{Erstes Potenzgesetz: Gleiche Basis}{fundamentals_power_law_same_base}
  Potenzen mit gleicher Basis werden multipliziert beziehungsweise dividiert,
  indem die Exponenten addiert beziehungsweise subtrahiert werden:
  \begin{align*}
    a^m\cdot a^n&=a^{m+n},\\
    \frac{a^m}{a^n}&=a^{m-n}\qquad(a\neq0).
  \end{align*}
\end{TheoremBox}

\begin{ExampleBox}{Zum ersten Potenzgesetz}{fundamentals_power_law_same_base_example}
  \begin{enumerate}[label=(\alph*),itemsep=0.35em]
    \item $3^4\cdot3^2=3^{4+2}=3^6$
    \item $\displaystyle\frac{5^7}{5^3}=5^{7-3}=5^4$
    \item $\displaystyle\frac{2^3}{2^5}=2^{3-5}=2^{-2}$
  \end{enumerate}
\end{ExampleBox}

Als unmittelbare Folgerung erhält man für $a\neq0$ und $n\geq1$
\[
  a^{-n}=\frac{1}{a^n}.
\]
Ein negativer Exponent steht somit für den Kehrwert der Potenz mit positivem
Exponenten.

\begin{TheoremBox}{Zweites Potenzgesetz: Gleicher Exponent}{fundamentals_power_law_same_exponent}
  Potenzen mit gleichem Exponenten werden multipliziert beziehungsweise
  dividiert, indem die Basen zusammengefasst werden:
  \begin{align*}
    a^n\cdot b^n&=(ab)^n,\\
    \frac{a^n}{b^n}&=\left(\frac ab\right)^n\qquad(b\neq0).
  \end{align*}
\end{TheoremBox}

\begin{ExampleBox}{Zum zweiten Potenzgesetz}{fundamentals_power_law_same_exponent_example}
  \begin{enumerate}[label=(\alph*),itemsep=0.35em]
    \item $2^3\cdot5^3=(2\cdot5)^3=10^3$
    \item $\displaystyle\frac{6^4}{3^4}=\left(\frac63\right)^4=2^4$
  \end{enumerate}
\end{ExampleBox}

\begin{TheoremBox}{Drittes Potenzgesetz: Potenz einer Potenz}{fundamentals_power_law_power_of_power}
  Wird eine Potenz nochmals potenziert, werden die Exponenten multipliziert:
  \[
    (a^m)^n=a^{mn}=(a^n)^m.
  \]
\end{TheoremBox}

\begin{ExampleBox}{Zum dritten Potenzgesetz}{fundamentals_power_law_power_of_power_example}
  \begin{enumerate}[label=(\alph*),itemsep=0.35em]
    \item $(2^3)^4=2^{3\cdot4}=2^{12}$
    \item $(5^2)^5=5^{2\cdot5}=5^{10}$
  \end{enumerate}
\end{ExampleBox}

\begin{RemarkBox}
  Beim dritten Potenzgesetz entscheidet die Klammer über die Reihenfolge:
  \[
    (a^m)^n=a^{mn},
    \qquad
    a^{m^n}\text{ ist im Allgemeinen etwas anderes.}
  \]
\end{RemarkBox}

\subsection{Wurzeln}
\label{subsec:fundamentals_roots}

\begin{DefinitionBox}{$n$-te Wurzel}{fundamentals_nth_root}
  Für $n\in\N$ mit $n\geq2$ ist die $n$-te Wurzel aus $a$ die reelle Zahl
  $b$ mit
  \[
    \sqrt[n]{a}=b
    \quad\Longleftrightarrow\quad
    b^n=a
    \quad\text{und}\quad
    \bigl(n\text{ ungerade oder }b\geq0\bigr).
  \]
  gilt. Dabei heißt $a$ \emph{Radikand} und $n$ \emph{Wurzelexponent}. Für
  $n=2$ schreiben wir kurz $\sqrt a$ statt $\sqrt[2]{a}$.

  \begin{center}
    \begin{tikzpicture}[font=\small,>=stealth]
  \node (root) at (0,0) {\Huge $\sqrt[n]{a}$};
  \node[anchor=east,text=THMBlue] (index) at (-1.45,0.95)
    {\Lang{Wurzelexponent}{root index}};
  \node[anchor=west,text=THMBlue] (radicand) at (1.55,-0.9)
    {\Lang{Radikand}{radicand}};
  \draw[->,THMBlue,line width=0.7pt] (index.south east) -- (-0.42,0.28);
  \draw[->,THMBlue,line width=0.7pt] (radicand.north west) -- (0.48,-0.12);
\end{tikzpicture}

  \end{center}
\end{DefinitionBox}

Ob eine reelle Wurzel existiert, hängt bei negativen Radikanden vom
Wurzelexponenten ab:
\[
  \begin{array}{c@{\qquad}c}
    n\text{ gerade}: & \sqrt[n]{a}\text{ ist nur für }a\geq0\text{ definiert},\\[0.4em]
    n\text{ ungerade}: & \sqrt[n]{a}\text{ ist für alle }a\in\R\text{ definiert}.
  \end{array}
\]

\begin{ExampleBox}{Gerade und ungerade Wurzeln}{fundamentals_even_odd_roots}
  \begin{enumerate}[label=(\alph*),itemsep=0.35em]
    \item $\sqrt{9}=3$, denn $3^2=9$.
    \item $\sqrt{-9}$ existiert nicht in $\R$.
    \item $\sqrt[3]{-8}=-2$, denn $(-2)^3=-8$.
  \end{enumerate}
\end{ExampleBox}

\subsection{Wurzeln als Potenzen}
\label{subsec:fundamentals_roots_as_powers}

Die Schreibweise für rationale Exponenten soll mit dem dritten Potenzgesetz
verträglich sein. Für $a>0$ suchen wir daher einen Exponenten $r$, für den
$(a^r)^n=a$ gilt. Nach dem dritten Potenzgesetz ist
\[
  (a^r)^n=a^{rn}=a^1
  \quad\Longrightarrow\quad
  rn=1
  \quad\Longrightarrow\quad
  r=\frac1n.
\]
Somit ist $a^{1/n}$ genau die Zahl, deren $n$-te Potenz $a$ ergibt, also die
$n$-te Wurzel aus $a$.

\begin{KeyPointBox}{Rationale Exponenten}{fundamentals_rational_exponents}
  Für $a>0$, $m\in\Z$ und $n\in\N$ mit $n\geq2$ gilt
  \[
    \sqrt[n]{a}=a^{1/n},
    \qquad
    a^{m/n}=\sqrt[n]{a^m}=\left(\sqrt[n]{a}\right)^m.
  \]
\end{KeyPointBox}

Damit folgen die Wurzelgesetze direkt aus den Potenzgesetzen:
\begin{align*}
  \sqrt[n]{a}\,\sqrt[n]{b}&=\sqrt[n]{ab},
  &
  \frac{\sqrt[n]{a}}{\sqrt[n]{b}}&=\sqrt[n]{\frac ab}\quad(b\neq0),\\
  \sqrt[n]{\sqrt[m]{a}}&=\sqrt[nm]{a}.&&
\end{align*}

\begin{RemarkBox}
  Wurzeln dürfen über Produkte und Quotienten verteilt werden, aber nicht über
  Summen oder Differenzen. Im Allgemeinen gilt
  \[
    \sqrt{a+b}\neq\sqrt a+\sqrt b.
  \]
\end{RemarkBox}

\begin{ExampleBox}{Rationale Exponenten}{fundamentals_rational_exponents_example}
  \begin{enumerate}[label=(\alph*),itemsep=0.35em]
    \item $27^{2/3}=\left(\sqrt[3]{27}\right)^2=3^2=9$
    \item $9^{-1/2}=\displaystyle\frac1{\sqrt9}=\frac13$
    \item $\sqrt{\sqrt[3]{64}}=\sqrt[6]{64}=2$
  \end{enumerate}
\end{ExampleBox}
